\section*{第 \ref{ch:b3:point-groups} 章——对称性与点群}

\begin{solution}{exo:b3:point-groups:1}
乙烯：三条相互垂直的 $C_2$ 轴、三个镜面、$i$：$D_{2h}$。\ce{XeF4}（正方形
平面）：$C_4$，四个 $C_2 \perp C_4$，$\sigma_h$，$2\sigma_v$，$2\sigma_d$，$i$，$S_4$：
$D_{4h}$。\ce{PCl5}（三角双锥）：$C_3$，三个 $C_2$，$\sigma_h$，$3\sigma_v$，
$S_3$：$D_{3h}$。1,3,5-三氯苯：对称元素与 \ce{BF3} 相同：$D_{3h}$。
\end{solution}

\begin{solution}{exo:b3:point-groups:2}
交叉式构象的乙烷：$C_3$，三个与主轴垂直的 $C_2$，三个 $\sigma_d$，$i$，$S_6$：
$D_{3d}$。重叠式构象：$C_3$，三个 $C_2$，$\sigma_h$，三个 $\sigma_v$：$D_{3h}$。
椅式构象的环己烷：$D_{3d}$。
\end{solution}

\begin{solution}{exo:b3:point-groups:3}
\ce{SO3}（$D_{3h}$）：否。\ce{SO2}（$C_{2v}$）：是。\ce{PCl3}（$C_{3v}$）：是。
\ce{XeF4}（$D_{4h}$）：否。\ce{CH2Cl2}（$C_{2v}$）：是。顺-1,2-二氯乙烯
（$C_{2v}$）：是；反式（$C_{2h}$）：否。
\end{solution}

\begin{solution}{exo:b3:point-groups:4}
$C_2(z) = \mathrm{diag}(-1,-1,1)$, $\chi = -1$; $i = \mathrm{diag}(-1,-1,-1)$,
$\chi = -3$; $\sigma_h(xy) = \mathrm{diag}(1,1,-1)$, $\chi = 1$.
\end{solution}

\begin{solution}{exo:b3:point-groups:5}
利用上题中的对角矩阵：$C_2i = \sigma_h$，$C_2\sigma_h =
i$，$i\sigma_h = C_2$，每个操作都是自己的逆，且任意两个操作的乘积
可交换。该群为交换群，故 $X^{-1}AX = A$ 对所有 $X$ 成立：每个
操作各自单独构成一个类（共四个类，四个
一维不可约表示）。
\end{solution}

\begin{solution}{exo:b3:point-groups:6}
取 $\sigma_v(1)$ 为 $xz$ 平面。点 $(1,0)$ 经 $C_3$ 变到 $(-\frac12,
\frac{\sqrt3}2)$，再经 $\sigma_v(1)$ 变到 $(-\frac12,-\frac{\sqrt3}2)$；另一顺序下
则先变到 $(1,0)$，再变到 $(-\frac12,\frac{\sqrt3}2)$。两个乘积是
两个不同镜面（$\sigma_v(3)$ 与 $\sigma_v(2)$）的反射：该群
不可交换，这正是其反射构成一个含三个元素的类的原因。
\end{solution}

\begin{solution}{exo:b3:point-groups:7}
扭转联苯保留沿环间键的 $C_2$ 以及两个
与之垂直的 $C_2$，它们平分两环平面之间的夹角；所有
镜面以及平面式或垂直式构型的反演中心都
已丧失。若有三个相互垂直的 $C_2$ 而无象转轴，则该群为
$D_2$：该分子具有手性（其两种扭转构型互为对映异构体，
当邻位大体积取代基阻碍旋转时可以拆分）。
\end{solution}

\begin{solution}{exo:b3:point-groups:8}
$d_{z^2}$ 与 $d_{x^2-y^2}$：$a_1$（$x^2$、$y^2$、$z^2$ 属 $A_1$）；$d_{xy}$：$a_2$；$d_{xz}$：
$b_1$；$d_{yz}$：$b_2$。
\end{solution}

\begin{solution}{exo:b3:point-groups:9}
$1 + 1 + 4 + 9 + 9 = 24 = h$. $\sum g\,\chi_{A_1}\chi_{T_2} = 1\cdot3 + 8\cdot0 + 3(-1)
+ 6(-1) + 6(1) = 0$.
\end{solution}

\begin{solution}{exo:b3:point-groups:10}
在与交线垂直的平面内，关于一条直线的反射
（直线倾角为 $\alpha$）将极角 $\phi$ 映为 $2\alpha - \phi$。两次反射，倾角依次为
$\alpha$ 与 $\alpha + \theta$：$\phi \to 2\alpha - \phi \to 2(\alpha + \theta) - (2\alpha - \phi) =
\phi + 2\theta$，即绕该直线的旋转 $2\theta$。两个夹角为
$60^\circ$ 的竖直镜面因此生成一个转角为 $120^\circ$ 的旋转，即 $C_3$。
\end{solution}

\begin{solution}{exo:b3:point-groups:11}
C=C=C 轴为 $C_2$，也为 $S_4$（转过 $90^\circ$，交换
两个 \ce{CH2} 平面，再经中心碳原子所在平面反射）；两个
\ce{CH2} 平面为 $\sigma_d$；两个与轴垂直的 $C_2$ 平分它们。
阶为 8：$D_{2d}$，非手性的。在 \ce{ClHC=C=CHCl} 中，镜面、$S_4$ 与轴向
$C_2$ 均因取代而被破坏；仅剩与轴垂直的
一个 $C_2$ 保留，群为 $C_2$，具有手性——轴手性丙二烯。
\end{solution}

\begin{solution}{exo:b3:point-groups:12}
\ce{SF6}：$O_h$（$h = 48$）。\ce{SF5Cl}：过 Cl 的 $C_4$ 及四个 $\sigma_v$：
$C_{4v}$（$h = 8$）。反式-\ce{SF4Cl2}：$D_{4h}$（$h = 16$）。顺式-\ce{SF4Cl2}：$C_{2v}$（$h
= 4$）。每个点群都保留 $O_h$ 的操作的一个子集，并对乘积封闭，
且 8、16、4 都能整除 48。有极性的：\ce{SF5Cl} 与顺式-\ce{SF4Cl2}。
\end{solution}

\begin{solution}{pb:b3:point-groups:1}
\textbf{1.} 沿过 C 与各 H 的四条体对角线；每条给出
$C_3$ 与 $C_3^2$：共 8 个旋转。
\textbf{2.} 沿过面心的 $x$、$y$、$z$ 轴：$C_2(z)$ 将 $(1,1,1)
\to (-1,-1,1)$，为一个氢。以 $90^\circ$ 绕 $z$ 轴旋转，再
关于 $xy$ 面的反射将 $(1,1,1) \to (-1,1,1) \to (-1,1,-1)$，为一个氢：一个
$S_4$；连同 $S_4^3$，共 6 个操作。
\textbf{3.} 平面 $x = \pm y$、$y = \pm z$、$x = \pm z$；$x = y$ 包含
$(1,1,1)$ 与 $(-1,-1,1)$：每个平面各包含两个 C--H 键。
\textbf{4.} $1 + 8 + 3 + 6 + 6 = 24$.
\textbf{5.} 反演将 $(1,1,1)$ 映为 $(-1,-1,-1)$，它不是氢原子：
无 $i$。
\textbf{6.} $E$; $8C_3$; $3C_2$; $6S_4$; $6\sigma_d$.
\textbf{7.} \ce{CH3Cl}：$C_{3v}$，$h = 6$。
\textbf{8.} \ce{CH2Cl2}：$C_{2v}$，$h = 4$。
\textbf{9.} \ce{CHCl3}：$C_{3v}$；\ce{CCl4}：$T_d$，$h = 24$。
\textbf{10.} \ce{CH2FCl}：$C_s$（$h = 2$）；\ce{CHFClBr}：$C_1$（$h = 1$）。
\textbf{11.} 各自保留 $T_d$ 中那些与
取代方式相容的操作；6、4、6、24、2、1 均整除 24。
\textbf{12.} 三个氢原子，被 $C_3$、$C_3^2$ 及三个
$\sigma_v$ 所置换。
\textbf{13.} 除 \ce{CCl4} 外：\ce{CH3Cl} 与 \ce{CHCl3} 沿 $C_3$ 轴，
\ce{CH2Cl2} 沿 $C_2$ 轴，\ce{CH2FCl} 在镜面内，
\ce{CHFClBr} 则无指定方向。
\textbf{14.} 仅 \ce{CHFClBr}（$C_1$，无 $S_n$）。
\textbf{15.} 两个对映异构体。
\textbf{16.} 其两个氢等价：过 C、F 与 Cl 的平面
平分 H--C--H 键角，是一个镜面。
\textbf{17.} 否：四个互不相同的取代基只保留 $E$；无 $S_n$ 时该
分子具有手性（且可能有极性）。
\textbf{18.} $(x, y, z)$：$T_2$。
\textbf{19.} 两个氢原子是四面体某条棱的两端；24
个操作可将任一条棱映到任一条其他棱（6 条棱构成一组）：只有
一种 \ce{CH2Cl2}。
\textbf{20.} \ce{CH2FCl}：一种。\ce{CHFClBr}：两种，为一对对映异构体。
\textbf{21.} 三种：邻位（$C_{2v}$）、间位（$C_{2v}$）、对位（$D_{2h}$）。
\textbf{22.} 三种：1,2,3-（$C_{2v}$）、1,2,4-（$C_s$）、1,3,5-（$D_{3h}$）。
\textbf{23.} 两种（两个 Cl 原子相邻或相对）。已知
仅有一种 \ce{CH2Cl2}，这排除了平面碳，并支持了
于 1874 年提出的四面体碳。
\textbf{24.} $1^2 + 1^2 + 2^2 + 3^2 + 3^2 = 24$：\textbf{$T_d$ 的阶，即甲烷的
点群，$h = 24$}。
\end{solution}
